\documentclass[11pt,a4paper]{article}
\usepackage[utf8]{inputenc}
\usepackage{geometry}
\geometry{margin=1in}
\usepackage{hyperref}
\usepackage{listings}
\usepackage{xcolor}
\usepackage{amsmath}
\usepackage{enumitem}
\usepackage{tocloft}

\title{\textbf{PORT-HUB.COM} \\ Universal Adaptive Port Hub}
\author{Atrographer}
\date{\today}

\definecolor{codegreen}{rgb}{0,0.6,0}
\definecolor{codegray}{rgb}{0.5,0.5,0.5}
\definecolor{codepurple}{rgb}{0.58,0,0.82}
\definecolor{backcolour}{rgb}{0.95,0.95,0.92}

\lstdefinestyle{mystyle}{
    backgroundcolor=\color{backcolour},  
    commentstyle=\color{codegreen},
    keywordstyle=\color{magenta},
    numberstyle=\tiny\color{codegray},
    stringstyle=\color{codepurple},
    basicstyle=\ttfamily\footnotesize,
    breakatwhitespace=false,        
    breaklines=true,                
    captionpos=b,                   
    keepspaces=true,                
    numbers=left,                   
    numbersep=5pt,                 
    showspaces=false,               
    showstringspaces=false,
    showtabs=false,                 
    tabsize=2
}

\lstset{style=mystyle}

\begin{document}

\maketitle

\tableofcontents

\begin{abstract}
PORT-HUB.COM is a port/adapter system. It discovers, types, registers, and interconnects components in a project — turning Python objects, modules, and project files (Vite + React + Render.com stacks) into ports with connections.
\end{abstract}

\section{Quick Start}

\begin{lstlisting}[language=Python]
from port_type_hub import PortHub

PortHub.auto_discover()

PortHub.auto_hub_files(".")

PortHub.status()
\end{lstlisting}

\section{Core Thesis}

PORT-HUB turns project artifacts into a graph of typed ports. It bridges frontend files (JSX/TSX files, configs) with Python orchestration through:

\begin{itemize}
    \item Port: Any object, module, or file (via FilePortWrapper).
    \item PortType: Typing system (vite:config, react:component, render:service).
    \item Dynamic Connection: Runtime method injection so ports connect.
    \item Auto-Discovery + Auto-Connect: For web stacks.
\end{itemize}

\section{auto\_hub\_files()}

\begin{lstlisting}[language=Python]
PortHub.auto_hub_files(
    root_dir=".",
    extensions=None,
    max_depth=10
)
\end{lstlisting}

\subsection{What It Does}
\begin{enumerate}
    \item Scans the project directory.
    \item Skips node\_modules, \_\_pycache\_\_, .git, etc.
    \item Registers files as FilePortWrapper objects (with .read(), path metadata).
    \item Applies type inference from default\_port\_types.py.
    \item Calls \_auto\_connect\_compatible() to wire ports.
\end{enumerate}

\subsection{Supported Types}

\begin{itemize}
    \item vite.config.js/ts \( \longrightarrow \) vite:config
    \item *.jsx / *.tsx \( \longrightarrow \) react:component
    \item Files with render, deploy \( \longrightarrow \) render:service
    \item *.vue / *.svelte \( \longrightarrow \) framework components
    \item index.html \( \longrightarrow \) web:static
    \item Python files \( \longrightarrow \) python:module
\end{itemize}

\section{Dynamic Connections}

After auto-hubbing, ports connect:

\begin{lstlisting}[language=Python]
for v in vite_ports:
    for r in react_ports:
        PortHub.connect(v, r, bidirectional=True)
    for ren in render_ports:
        PortHub.connect(v, ren, bidirectional=True)
\end{lstlisting}

connect() injects proxy methods:

- vite\_port.to\_react-component\_build()
- react\_port.to\_render-service\_deploy()

\section{Example in a Vite + React Project}

\begin{lstlisting}[language=Python]
from port_type_hub import PortHub

def main():
    hub = PortHub()
   
    hub.auto_hub_files(".")
   
    vite_config = hub.get_port_by_type("vite:config")
    react_components = hub.get_ports_by_type("react:component")
   
    hub.status()
   
if __name__ == "__main__":
    main()
\end{lstlisting}

\section{Advanced Usage}

\begin{itemize}
    \item Manual registration: PortHub.register\_port(name, obj, port\_type)
    \item Custom filters with ModFilter
    \item Extend types in default\_port\_types.py
    \item Use in build scripts and orchestration
\end{itemize}

\section{Project Structure}

Typical Vite + React + Render project:

\begin{verbatim}
my-vite-react-app/
├── index.html
├── package.json
├── vite.config.js / .ts
├── tsconfig.json
├── postcss.config.js
├── tailwind.config.js
├── .env
├── public/
│   └── vite.svg
└── src/
    ├── main.tsx
    ├── App.tsx
    ├── components/
    ├── pages/
    ├── assets/
    ├── utils/
    ├── hooks/
    └── styles/
\end{verbatim}

\section{Usage}

After using Port Hub the project components are connected.

\end{document} 
